% LSPV-4 REPORT::REACQ1::2026-10-01T15:46:00-07:00
% !TEX program = xelatex
\documentclass[11pt]{article}
\usepackage[margin=0.88in]{geometry}
\usepackage{amsmath,amssymb}
\begin{document}
\section*{LSPV-4: deterministic reacquisition attack on FP/FNP Search-G and RCCN}
\textbf{STATUS: OPEN.} P versus NP remains open. No unrestricted SAT lower bound is proved here.

\paragraph{Predecessor preflight.}
The newest LSPV-3 report, \texttt{NOGUARD1}, is substantive, has the required header, and ends with the required final verification sentinel. Predecessor recovery is therefore not requested.

\paragraph{CORPUS REVIEWED.}
I surveyed the repository tree and governing files \texttt{AGENTS.md}, \texttt{CODEX.md}, and \texttt{.codex/skills/pnp-research/SKILL.md}; root \texttt{main.tex}; current TASK-001 through TASK-007 surfaces; newest/relevant MAIN snapshots; materially relevant LSPV-1--5 material including PAIR2COVER1, WEDGE1, NOGUARD1, ACQ1, REREAD1, PCLOSE1-R1, and RANDCLOSE1; governance; HEC/VSC/G-PPC/CRS-PPC/CCSPC/RCCN; the FP/FNP parallel proof spine and Search-G precommitment spine; and repository searches for STORE/REACQUIRE, Search-G, FP/FNP, RCCN, current MAIN, and duplicate causal claims. The verified recovery PCLOSE1-R1 and newer RANDCLOSE1 are retained.

\paragraph{Two obstacles and literature preflight.}
The two obstacles attacked here are (i) whether the exact RCCN STORE-or-REACQUIRE dichotomy can be promoted to a superpolynomial theorem for a deterministic FP selector, and (ii) whether Search-G future-bit precommitment prevents polynomial reacquisition. Primary literature checked includes Bellare--Goldwasser on the distinction between generic decision/search reductions, Valiant--Vazirani on search/uniqueness phenomena, and Aaronson--Wigderson on algebrization. The Bellare--Goldwasser caveat does not rescue Search-G under the explicit counterfactual: for a polynomially balanced FNP relation, the prefix-extension language is in NP, so under $P=NP$ a witness can be recovered in polynomial time by adaptive prefix queries. Recent search-to-decision work likewise treats witness extraction through adaptive NP queries rather than providing a static lower bound.

\paragraph{PROVED/EXACT: same-selector reacquisition theorem.}
Let $R(x,y)$ be a polynomially balanced FNP relation and suppose a deterministic machine $S$ solves it in time
\[
T_S(n)\le n^c.
\]
For every $x$ in the search domain, write $y=S(x)$. Consider any later audit challenge $g$ for which the demanded answer has the form
\[
A(x,y,g)
\]
with $A$ deterministic polynomial time. If the original static input $x$ remains readable, then after discarding $y$ completely a legal implementation can answer the audit by rerunning $S(x)$ and evaluating $A(x,S(x),g)$. Its time is
\[
\boxed{T_S(|x|)+T_A(|x|+|S(x)|+|g|)=\operatorname{poly}(|x|+|g|).}
\]
Because $S$ is deterministic, the reacquired witness is the same witness as before the cut. Thus SAME\_REALIZATION at the output level does not force persistent witness storage.

The statement extends to any polynomial number $q(n)$ of delayed audits: rerun $S$ independently for each audit, for total time at most
\[
\boxed{q(n)\,T_S(n)+\sum_{j\le q(n)}T_{A_j}(\operatorname{poly}(n))=\operatorname{poly}(n).}
\]
No memoization is required.

\paragraph{PROVED/EXACT: deterministic trace-coordinate reacquisition.}
Let the deterministic computation of $S(x)$ have configurations
\[
C_0(x),C_1(x),\ldots,C_{T_S(n)}(x).
\]
Given a later index $i\le T_S(n)$, a static machine can reproduce $C_i(x)$ by rerunning $S$ from $C_0(x)$ for $i$ steps. Hence any polynomial number of requested coordinates or polynomial-time predicates of individually addressable configurations can be serviced with polynomial total recomputation. If an audit explicitly requests the whole polynomial-length trace, the trace itself can be regenerated in polynomial time and has only polynomial output length. Therefore moving the delayed audit from the witness $y$ to a deterministic polynomial-time computation record does not by itself yield a superpolynomial reacquisition cost.

\paragraph{Consequence for RCCN-6 and Search-G.}
RCCN-6's STORE-or-REACQUIRE theorem remains \textbf{PROVED/EXACT in its declared model}: after a cut, enough information must be stored or obtained again. The new hostile result concerns the next arrow:
\[
\boxed{\text{REACQUIRE is necessary}\Longrightarrow\text{REACQUIRE is superpolynomial}.}
\]
For deterministic polynomial-time selectors on static readable inputs, that generic arrow is \textbf{FALSIFIED/RETIRED}. Re-execution is already a polynomial reacquisition strategy.

Accordingly, the Search-G chain
\[
P=NP\Rightarrow FP=FNP\Rightarrow S\in FP\Rightarrow
\text{future }G\Rightarrow\text{STORE or REACQUIRE}
\]
remains valid through the dichotomy, but the proposed promotion to a superpolynomial resource cannot follow merely from future-bit precommitment, SAME\_REALIZATION of $S(x)$, or delayed audit. A surviving theorem must prohibit or quantitatively defeat ordinary rerun without imposing a restriction absent from arbitrary static SAT algorithms.

\paragraph{Counterfactual closure.}
Under $P=NP$, every polynomially balanced FNP relation has a deterministic polynomial-time selector by the prefix-extension construction. Therefore the all-witness search spine does not escape PCLOSE1-R1 merely by replacing a decision bit with a polynomial witness. If a construction requires an arbitrary selector's actual output, deterministic rerun reproduces that selector's output. If it requires merely some valid witness, the $FP=FNP$ selector reacquires one in polynomial time. Canonicalization does not help because canonical SAT witnesses are also obtainable by polynomially many prefix decisions under the counterfactual.

\paragraph{Guards and limits.}
This theorem does not show that every conceivable reacquisition task is easy. It does not cover a demanded object of superpolynomial explicit length, a superpolynomial realized history, an audit whose own evaluation is independently proved superpolynomial, or an input that is genuinely unavailable after the cut in the target computation model. But each escape requires a new model-preserving bridge. Declaring the static SAT input unreadable after a cut would contradict the arbitrary-static-algorithm target; making the history succinctly represent superpolynomial work risks a succinct-input/NEXP change of problem.

\paragraph{Hostile controls.}
DEAD/COPY/DELAY only inflate the record and remain rerunnable. INTERP/RECOMP explicitly implement the bypass. Static reread preserves $x$. Memoization can only improve the reacquisition cost. Group compression and pointer jumping can shorten special histories further. Gaussian/Tseitin, Horn, 2-SAT, and bounded-treewidth controls show that long syntactic dependency or audit chains can coexist with globally easy computation. Stagewise nondeterminism is no stronger here: under $P=NP$, the corresponding polynomially balanced search can be deterministically selected in polynomial time. Padding cannot convert the polynomial rerun into a superpolynomial lower bound measured against the padded input.

\paragraph{PQ and cross-domain classification.}
\textbf{EXACT OBJECT/THEOREM TRANSFER:} fixed-code Schur products, quadratic relation kernels, typed puncturing/shortening, POLAR1/WEDGE1 algebra, collision Cauchy--Schwarz, and the deterministic rerun/reacquisition theorem above.
\textbf{ANALOGY/CANDIDATE:} syzygies/Betti invariants as succinct separators, special-soundness collision extraction, LWE-style external extraction, communication pointer chasing, cell-probe lower bounds, and proof-search lower bounds. Each needs a model-preserving theorem.
\textbf{FAILED TRANSFER:} cryptographic hardness/random-oracle assumptions, VLSI/distributed locality, online no-reread state bounds, and proof-system-specific hardness as unrestricted static SAT runtime lower bounds.

\paragraph{SU/vMS, acquisition, and reconstruction retained.}
PAIR2COVER1/WEDGE1 sharpen fixed-Hadamard algebra but do not source-transcribe the exact SU09/vMS invocation or close the stacked Schur rank. ACQ1 remains CONDITIONAL on the exact legal-state evaluator and SAT compilation exponent. Sparse-$Q$ normalization remains OPEN. NOGUARD1 retains the dependency-sensitive reconstruction ledger and the trivial $A\circ D_j$ NO-evaluator baseline. None is strengthened by the present causal falsification.

\paragraph{Relativization and algebrization audit.}
REACQ1 is explicitly compatible with semantic-oracle replacement and therefore is not a nonrelativizing hinge. The active Schur/polar machinery remains algebraic. Aaronson--Wigderson's algebrization barrier remains live. Any final contradiction still needs an arrow that fails under the relevant semantic/algebraic oracle strengthening while holding for every ordinary polynomial-time SAT algorithm.

\paragraph{DUPLICATION / CONTRADICTION AUDIT.}
REREAD1 already showed that causal pointer dependence does not force persistent state; PCLOSE1-R1 proved deterministic polynomial-depth closure under $P=NP$; RANDCLOSE1 extended closure to bounded-error randomized polynomial-depth computations with adaptive NP calls. The FP/FNP spine itself already states that the missing theorem is a lower bound on reacquisition. REACQ1 does not duplicate those results: it gives the explicit same-selector rerun simulation and trace-coordinate reacquisition that directly falsify a generic superpolynomial REACQUIRE arrow. No contradiction with RCCN-6 exists because RCCN-6 states necessity of storage or reacquisition, not a superpolynomial lower bound on reacquisition.

\paragraph{STALE OR OVERSTATED MAIN CLAIMS; adopted/rejected/retired.}
Root MAIN's Search-G candidate language is acceptable only if its final STORE/REACQUIRE-to-superpolynomial arrow remains visibly OPEN. Any stronger reading that future-bit precommitment or SAME\_REALIZATION blocks deterministic rerun is stale. Adopted: exact RCCN dichotomy, FP/FNP equivalence under the declared convention, fixed-Hadamard algebra, charged ACQ1, and reconstruction guards. Rejected as insufficient: witness materialization, precommitment, delayed challenge, deterministic trace provenance, and all-witness wording as standalone hardness. FALSIFIED/RETIRED: generic promotion of REACQUIRE necessity to superpolynomial cost for polynomial-time deterministic selectors on static readable inputs.

\paragraph{P=NP-compatible components without claiming P=NP.}
Polynomial FNP witness selection under the counterfactual, same-selector rerun, trace regeneration, static reread, canonical self-reduction, polynomial adaptive SAT calls, bounded-error polynomial causal cascades, stagewise witness selection, memoization, group compression, and ACQ1 NP search are all compatible with hypothetical $P=NP$. This is not evidence that $P=NP$.

\paragraph{Strict status and next falsification target.}
\textbf{PROVED/EXACT:} same-selector witness reacquisition and deterministic trace-coordinate reacquisition with the stated polynomial costs; RCCN STORE-or-REACQUIRE remains exact in its own model.
\textbf{KNOWN THEOREM:} under the standard polynomially balanced search convention, $P=NP$ implies polynomial-time FNP witness search by prefix-extension.
\textbf{CONDITIONAL:} ACQ1's charged acquisition ledger and the current reconstruction exponent accounting.
\textbf{CANDIDATE:} a Search-G audit whose reacquisition cost survives deterministic rerun and static reread by an independently established model-preserving theorem.
\textbf{FALSIFIED/RETIRED:} generic STORE-or-REACQUIRE necessity as a route to superpolynomial cost for an FP selector; future-bit precommitment as a standalone blocker of rerun.
\textbf{OPEN:} exact SU09/vMS invocation, stacked Schur rank/common kernel, sparse-$Q$ normalization, ACQ1 compilation exponent, reconstruction improvement beyond $A\circ D_j$, arbitrary-algorithm transfer, Search-G's superpolynomial resource arrow, and a nonrelativizing/nonalgebrizing hinge. Hence P versus NP remains open.

The next falsification target is the strongest proposed Search-G audit after explicitly granting the adversary a \emph{rerun oracle for its own deterministic selector}. If the audit can still force superpolynomial total cost with the original static input readable and all rerun time fully charged, state the exact theorem and resource. If it cannot, retire Search-G's causal-resource route rather than adding further precommitment terminology.

\begin{thebibliography}{5}
\bibitem{BG94} M. Bellare and S. Goldwasser, \emph{The Complexity of Decision Versus Search}, SIAM J. Comput. 23(1), 1994.
\bibitem{VV86} L. Valiant and V. Vazirani, \emph{NP is as Easy as Detecting Unique Solutions}, Theor. Comput. Sci. 47, 1986.
\bibitem{GK24} S. Gharibian and J. Kamminga, \emph{BQP, meet NP: Search-to-decision reductions and approximate counting}, 2024.
\bibitem{AW09} S. Aaronson and A. Wigderson, \emph{Algebrization: A New Barrier in Complexity Theory}, ACM Trans. Comput. Theory 1(1), 2009.
\bibitem{vMS25} D. van Melkebeek and N. Mocelin Sdroievski, \emph{Instance-Wise Hardness and Refutation versus Derandomization for Arthur--Merlin Protocols}, Computational Complexity 34:15, 2025.
\end{thebibliography}
\end{document}
% END_OF_VERIFIED_REPORT::LSPV-4::REACQ1